%% SVD and QR
%
% The two decompositions every algorithm here is made of, applied to one
% tensor of a chain: T with its left bond, its physical index and its right
% bond.
%
% Above, the SVD with the left bond and the physical index grouped:
% T = U S V^dagger. U keeps both of those and is isometric from the left; S is
% the singular values, the centre on a bond; V^dagger carries the right bond
% and is isometric from the right. Truncate S and this is the step that keeps
% a bond dimension bounded.
%
% Below, the QR of the same grouping: T = Q R. Q is isometric from the left
% and R is what is left over -- not diagonal, not isometric, a plain array --
% and it is what gets pushed into the next site when a chain is brought to
% canonical form. QR is the cheaper of the two when nothing is truncated.
%
% The arrows are not written: each bond takes its direction from the triangle
% it leaves, and an open bond of a triangle points in.
\documentclass[border=4pt]{standalone}
\usepackage{amsmath,amssymb}
\usepackage{tikz-tensors}
\input{notation}
\begin{document}
\begin{tikzpicture}
  \tnstack{T}{1}{ket}
  \tnlayer{T}{ket}{coef/$T$}
  \tnopen{left}{T}
  \tnopen{right}{T}
  \tnopen{down}{T}
  \tneq
  \tnstack{U}{3}{ket}
  \tnlayer{U}{ket}{canl/$U$, centrebond/$S$, canr/$V^{\dagger}$}
  \tnconnect[along=gauge]{U}
  \tnopen[edge=gauge]{left}{U-ket-1}
  \tnopen[edge=gauge]{right}{U-ket-3}
  \tnopen{down}{U-ket-1}
  \tnbreak
  \tnstack{P}{1}{ket}
  \tnlayer{P}{ket}{coef/$T$}
  \tnopen{left}{P}
  \tnopen{right}{P}
  \tnopen{down}{P}
  \tneq
  \tnstack{Q}{2}{ket}
  \tnlayer{Q}{ket}{canl/$Q$, coef/$R$}
  \tnconnect[along=gauge]{Q}
  \tnopen[edge=gauge]{left}{Q-ket-1}
  \tnopen{right}{Q-ket-2}
  \tnopen{down}{Q-ket-1}
\end{tikzpicture}
\end{document}
